\section{Практическая часть}

Лабораторная установка состоит из блока для изучения разрывных колебаний, генератора импульсов и двухканального осциллографа. Блок включает схему мультивибратора (триггера) и схему кипп-реле; регулятором входного напряжения $E$ задаётся режим мультивибратора или триггера.

\subsection{Режим мультивибратора}

С помощью регулятора 1 выбрали значение постоянного напряжения $E$ таким образом, чтобы кривая $\varphi(x)$ соответствовала режиму мультивибратора. На осциллографе наблюдали колебания напряжения и тока (\cref{fig:pr:1a}), а в режиме $X$--$Y$ --- фазовую плоскость мультивибратора (\cref{fig:pr:1b}).

\begin{figure}[H]
	\centering
	\begin{minipage}{0.45\textwidth}
		\centering
		\includegraphics[width=0.95\linewidth]{task_1}
		\subcaption{Осциллограммы напряжения и тока}
		\label{fig:pr:1a}
	\end{minipage}
	\begin{minipage}{0.45\textwidth}
		\centering
		\includegraphics[width=0.95\linewidth]{task_1_phase_plane}
		\subcaption{Фазовая плоскость}
		\label{fig:pr:1b}
	\end{minipage}
	\caption{Мультивибратор: осциллограммы напряжения и тока (а) и фазовая плоскость (б)}
	\label{fig:pr:1}
\end{figure}

По осциллограммам \cref{fig:pr:1} видно:
\[
U_{\max} = 1000~\text{мВ}, \qquad U_{\min} = 650~\text{мВ}, \qquad T = 170~\text{мкс}, \qquad I = ~\text{у.е.}
\]

Период колебаний тока совпадает с периодом колебаний напряжения.

\subsection{Режим триггера}

С помощью регулятора 1 перевели схему в режим триггера и подключили генератор прямоугольных импульсов ко входу мультивибратора. На вход подавали последовательность импульсов с частотой $f_\text{и} = 1$~кГц, то есть с периодом $T_\text{и} = 1$~мс, и амплитудой $A = 900$~мВ.

а) Длительность запускающего импульса задали равной 2~\% от периода повторения, $\tau_\text{и} = 20$~мкс. Полученная осциллограмма приведена на \cref{fig:pr:2}.

\begin{figure}[t]
	\centering
	\includegraphics[width=0.5\linewidth]{task_2_2pct_900mV}
	\caption{Осциллограммы входного сигнала (нижняя кривая) и напряжения на выходе триггера (верхняя кривая) при $\tau_\text{и} = 20$~мкс}
	\label{fig:pr:2}
\end{figure}

Триггер переключается на каждый запускающий импульс, а длительность импульса на выходе совпадает с периодом повторения входных импульсов:
\begin{align} \label{eq:pr:1}
	T_\text{тр} = T_\text{и} = 1~\text{мс}.
\end{align}

б--в) При уменьшении длительности запускающего импульса до 0,7~\% от периода ($\tau_\text{и} \approx 7$~мкс) переброс прекращается (\cref{fig:pr:3a}): за это время изображающая точка не успевает пройти участок медленных движений, необходимый для перехода в другое состояние равновесия, и после окончания импульса возвращается в исходное состояние.

При увеличении длительности до 3,7~\% изображающая точка остаётся в одном состоянии, но при подаче импульса наблюдается быстрое движение. Оно возникает из-за того, что за время импульса изображающая точка успевает дойти до точки $b$ (см. рис. 8 в методичке), после чего возвращается по траектории (I) в состояние 1.

\begin{figure}[t]
	\centering
	\begin{minipage}{0.45\textwidth}
		\centering
		\includegraphics[width=0.95\linewidth]{task_2_0p7_900mV}
		\subcaption{$\tau_\text{и} = 7$~мкс}
		\label{fig:pr:3a}
	\end{minipage}
	\begin{minipage}{0.45\textwidth}
		\centering
		\includegraphics[width=0.95\linewidth]{task_2_3p7_900mV}
		\subcaption{$\tau_\text{и} = 37$~мкс}
		\label{fig:pr:3b}
	\end{minipage}
	\caption{Влияние длительности запускающего импульса на работу триггера: при слишком коротком (а) и слишком длинном (б) импульсе схема перестаёт работать как триггер}
	\label{fig:pr:3}
\end{figure}

Таким образом, минимальная допустимая длительность запускающего импульса лежит в пределах 7--20~мкс, а максимальная --- 20--37~мкс.

Далее изменяли амплитуду запускающих импульсов при неизменной длительности $\tau_\text{и} = 20$~мкс. При $A = 500$~мВ переброса нет (\cref{fig:pr:4a}) --- амплитуды недостаточно, чтобы перевести изображающую точку в область притяжения другого состояния равновесия. При $A = 1{,}5$~В в выходном сигнале появляются срывы (\cref{fig:pr:4b}), и схема также перестаёт работать как триггер.

Пока схема работает как триггер, длительность и амплитуда выходного импульса от амплитуды запускающего не зависят: они определяются только параметрами самой схемы, а внешнее воздействие лишь меняет вид кривой $\varphi(x)$ и скорость движения изображающей точки на участках быстрых движений.

\begin{figure}[t]
	\centering
	\begin{minipage}{0.45\textwidth}
		\centering
		\includegraphics[width=0.95\linewidth]{task_2_2pct_0p5V}
		\subcaption{$A = 0{,}5$~В}
		\label{fig:pr:4a}
	\end{minipage}
	\begin{minipage}{0.45\textwidth}
		\centering
		\includegraphics[width=0.95\linewidth]{task_2_2pct_1p5V}
		\subcaption{$A = 1{,}5$~В}
		\label{fig:pr:4b}
	\end{minipage}
	\caption{Влияние амплитуды запускающего импульса на работу триггера при $\tau_\text{и} = 20$~мкс: амплитуда ниже порога (а) и выше порога (б)}
	\label{fig:pr:4}
\end{figure}

Пороговые значения амплитуды: минимальная амплитуда запускающего импульса лежит между 0,5 и 0,9~В, максимальная --- между 0,9 и 1,5~В.

г) Для наблюдения деления частоты увеличили частоту запускающих импульсов до $f_\text{и} = 11{,}8$~кГц ($T_\text{и} \approx 84{,}7$~мкс) и установили заполнение 11~\% ($\tau_\text{и} \approx 9{,}3$~мкс) при $A = 900$~мВ. В этом режиме переключение происходит на каждый второй запускающий импульс (\cref{fig:pr:5}), поэтому каждый уровень выходного сигнала держится в течение двух периодов входного:
\begin{align} \label{eq:pr:2}
	\tau_\text{вых} = 2T_\text{и} = 2 \cdot 84{,}7~\text{мкс} \approx 169~\text{мкс}.
\end{align}

\begin{figure}[t]
	\centering
	\includegraphics[width=0.5\linewidth]{task_2_division}
	\caption{Деление частоты на триггере при $f_\text{и} = 11{,}8$~кГц, $A = 900$~мВ и заполнении 11~\%}
	\label{fig:pr:5}
\end{figure}

\subsection{Режим кипп-реле}

Подготовили установку к изучению колебаний кипп-реле: подключили генератор ко входу схемы, а осциллограф --- к её выходу.

а) На вход подали сигнал с частотой $f_\text{и} = 24$~кГц ($T_\text{и} \approx 41{,}7$~мкс), амплитудой $A = 400$~мВ и заполнением 2~\% ($\tau_\text{и} \approx 0{,}83$~мкс). На выходе получили сигнал, приведённый на \cref{fig:pr:6}: короткий запускающий импульс формирует широкий прямоугольный импульс, длительность которого определяется только параметрами самой схемы и не зависит от длительности входного сигнала.

\begin{figure}[t]
	\centering
	\includegraphics[width=0.5\linewidth]{task_3_2pct}
	\caption{Сигнал на выходе кипп-реле при $f_\text{и} = 24$~кГц, $A = 400$~мВ и заполнении 2~\%}
	\label{fig:pr:6}
\end{figure}

б) Определили минимальную и максимальную длительность запускающего импульса, при которой схема работает как кипп-реле, то есть длительность выходного сигнала не зависит от длительности входного. При заполнении 0,19~\% ($\tau_\text{и} \approx 0{,}08$~мкс) схема формирует импульс с длительностью на полделения больше, следовательно, перестаёт функционировать как кипп-реле (\cref{fig:pr:7a}). При увеличении заполнения до 6,7~\% ($\tau_\text{и} \approx 2{,}8$~мкс) длительность выходного импульса уменьшается на полделения (\cref{fig:pr:7b}), схема перестаёт работать как кипп-реле.

\begin{figure}[t]
	\centering
	\begin{minipage}{0.45\textwidth}
		\centering
		\includegraphics[width=0.95\linewidth]{task_3_0p19pct}
		\subcaption{$\tau_\text{и} = 0{,}08$~мкс}
		\label{fig:pr:7a}
	\end{minipage}
	\begin{minipage}{0.45\textwidth}
		\centering
		\includegraphics[width=0.95\linewidth]{task_3_6p7pct}
		\subcaption{$\tau_\text{и} = 2{,}8$~мкс}
		\label{fig:pr:7b}
	\end{minipage}
	\caption{Работа кипп-реле при минимальной (а) и максимальной (б) длительности запускающего импульса}
	\label{fig:pr:7}
\end{figure}

Минимальная и максимальная длительность запускающего импульса:
\[
\tau_{\text{и}\,\min} \approx 0{,}08~\text{мкс}; \qquad \tau_{\text{и}\,\max} \approx 2{,}8~\text{мкс}.
\]

в) При уменьшении амплитуды запускающего импульса до 0,2~В спусковая схема не запускается: на выходе виден только короткий выброс, а широкий импульс не формируется (\cref{fig:pr:8a}). При $A = 0{,}21$~В схема запускается и работает как кипп-реле (\cref{fig:pr:8b}), поэтому минимальная амплитуда запускающего импульса
\[
A_{\min} = 0{,}21~\text{В}.
\]

\begin{figure}[t]
	\centering
	\begin{minipage}{0.45\textwidth}
		\centering
		\includegraphics[width=0.95\linewidth]{task_3_0p2V}
		\subcaption{$A = 0{,}20$~В}
		\label{fig:pr:8a}
	\end{minipage}
	\begin{minipage}{0.45\textwidth}
		\centering
		\includegraphics[width=0.95\linewidth]{task_3_0p21V}
		\subcaption{$A = 0{,}21$~В}
		\label{fig:pr:8b}
	\end{minipage}
	\caption{Порог запуска кипп-реле по амплитуде входного импульса}
	\label{fig:pr:8}
\end{figure}

При увеличении амплитуды входного импульса период выходного сигнала изменяется: при $A = 0{,}95$~В длительность соответствует начальной, а при $A = 0{,}99$~В период сигнала увеличивается на полделения (\cref{fig:pr:9}). При дальнейшем увеличении амплитуды схема перестаёт работать как кипп-реле, следовательно:
\[
A_{\max} \approx 0{,}95~\text{В}.
\]

\begin{figure}[t]
	\centering
	\begin{minipage}{0.45\textwidth}
		\centering
		\includegraphics[width=0.95\linewidth]{task_3_0p95V}
		\subcaption{$A = 0{,}95$~В}
		\label{fig:pr:9a}
	\end{minipage}
	\begin{minipage}{0.45\textwidth}
		\centering
		\includegraphics[width=0.95\linewidth]{task_3_0p99V}
		\subcaption{$A = 0{,}99$~В}
		\label{fig:pr:9b}
	\end{minipage}
	\caption{Искажение формы выходного импульса кипп-реле при амплитуде входного сигнала, близкой к максимальной}
	\label{fig:pr:9}
\end{figure}
